\documentclass[11pt,a4paper]{article}

\usepackage[margin=2.5cm]{geometry}
\usepackage{amsmath}
\usepackage{amssymb}
\usepackage{booktabs}
\usepackage{hyperref}
\usepackage{parskip}
\usepackage{microtype}
\usepackage{array}
\usepackage{cite}

\hypersetup{
    colorlinks=true,
    linkcolor=black,
    citecolor=black,
    urlcolor=blue
}

\title{\textbf{Dragonfly Algorithm}\\
       \large Synopsis: CITS4404 Part 1 Literature Review}
\author{James Wigfield \\ student ID: 23334375}
\date{April 2026}

\begin{document}

\maketitle

\noindent\textbf{Source paper:} S.\ Mirjalili, ``Dragonfly algorithm: a new meta-heuristic optimization technique for solving single-objective, discrete, and multi-objective problems,'' \textit{Neural Computing and Applications}, vol.~27, pp.~1053--1073, 2016.~\cite{mirjalili2016dragonfly}

\bigskip

\section*{1.\ What problem is the algorithm attempting to solve?}

The Dragonfly Algorithm (DA) is a swarm intelligence (SI) optimiser designed to
handle continuous, binary, and multi-objective optimisation problems.
Although algorithms like Particle Swarm Optimisation (PSO)~\cite{kennedy1995pso},
Ant Colony Optimisation (ACO), and Artificial Bee Colony (ABC) have been widely
studied, they each only model a narrow slice of the behaviours seen in natural swarms.
Dragonfly swarming in particular had not been explored computationally before this paper.
The No Free Lunch (NFL) theorem~\cite{wolpert1997nfl} also is also one of the motivations
for proposing this new algorithm, since no algorithm works best on every problem, a genuinely new approach
still has room to outperform existing ones on the right class of problems.

\bigskip



\section*{2.\ Why have previous attempts failed?}

Prior SI algorithms each capture some aspects of swarm behaviour, but none of them model all five survival behaviours at once.
For example, PSO only uses attraction towards a personal best and a global best (analogous to cohesion and food attraction), 
and while several improved PSO variants have tried to add separation, alignment, or cohesion operators independently~\cite{cui2009boid,kadrovach2002particle,cui2009alignment}, 
none of them really put all five behaviours into one method. Also, most existing algorithms don't clearly separate two different swarming modes: 
\textit{static} swarms that search in a small area (exploration) and \textit{dynamic} migratory swarms that move in one direction (exploitation).
PSO blurs this distinction with a single velocity update rule, making the exploration/exploitation trade-off hard to control directly.
Another gap is that most standard formulations have no concept of an \textit{enemy} (worst solution), 
so search agents are only pulled toward promising regions but never pushed away from bad ones, which can hurt diversity early in the search.

\bigskip

\section*{3.\ What is the new idea?}

The core idea is to model artificial dragonflies using \textbf{five behavioural operators} built on Reynolds' flocking rules~\cite{reynolds1987flocks}, with two extra terms added for food and enemy interactions:

\begin{enumerate}
    \item \textbf{Separation} ($S_i = -\sum_j (X - X_j)$): avoids collision with neighbours.
    \item \textbf{Alignment} ($A_i = \frac{1}{N}\sum_j V_j$): matches neighbours' velocities.
    \item \textbf{Cohesion} ($C_i = \frac{1}{N}\sum_j X_j - X$): moves toward the neighbourhood centre of mass.
    \item \textbf{Food attraction} ($F_i = X^+ - X$): steers toward the current best solution.
    \item \textbf{Enemy distraction} ($E_i = X^- + X$): steers away from the current worst solution.
\end{enumerate}

The \textbf{step vector} (analogous to PSO velocity) combines all five:
\[
    \Delta X_{t+1} = (s\,S_i + a\,A_i + c\,C_i + f\,F_i + e\,E_i) + w\,\Delta X_t
\]
where $w$ decays over iterations to shift from exploration to exploitation.
When a dragonfly has no neighbours it performs a \textbf{Lévy flight} ($X_{t+1} = X_t + \text{Levy}(d)\cdot X_t$), giving isolated agents a way to still explore the space.

Static swarming (high alignment weight $a$, low cohesion weight $c$) models exploration; dynamic swarming (low $a$, high $c$) models exploitation.
The neighbourhood radius increases over time, so what starts as many small independent sub-swarms gradually merges into one large convergent swarm.

Two variants are also presented.
The \textbf{Binary DA (BDA)} uses a v-shaped transfer function to convert the continuous step values into bit-flip probabilities --- v-shaped functions work better here because s-shaped ones tend to push particles toward extreme binary values~\cite{saremi2014transfer}.
The \textbf{Multi-Objective DA (MODA)} adds a Pareto archive and a grid system where food sources are picked from the least-crowded cells (to spread solutions out) and enemies from the most-crowded ones (to avoid revisiting dense areas), borrowing the archive approach from MOPSO~\cite{coello2004mopso}.

\bigskip

\section*{4.\ How is the new approach demonstrated?}

Each of the three algorithm variants is tested separately against established baselines.
DA is benchmarked on 19 standard test functions (unimodal, multi-modal, and composite) against PSO and GA, with the Wilcoxon rank-sum test used to check whether any differences are statistically significant.
BDA is tested on 13 of the same functions encoded in binary space, compared against Binary PSO and Binary GSA.
MODA is evaluated on five ZDT multi-objective benchmarks plus a real submarine propeller design problem, using IGD as the performance metric and comparing against MOPSO and NSGA-II.
The paper includes pseudocode, equations, and parameter settings in enough detail to replicate the experiments, and the source code is publicly available.

\bigskip

\section*{5.\ Results and validation}

\textbf{DA vs PSO vs GA (single-objective):} DA achieves the best or joint-best result on 13 of 19 functions, and Wilcoxon p-values confirm statistical significance over GA on virtually all functions ($p < 0.0002$).
Against PSO, differences are often not statistically significant on multi-modal functions (e.g.\ TF9, TF11, TF12), suggesting comparable exploration ability.
A notable exception is TF8 (Schwefel), where PSO substantially outperforms DA; this is the only case where a single-function weakness appears systematic rather than marginal.
On composite functions (TF14--TF19) the advantage over PSO is less consistent, which the author attributes honestly to the difficulty of these landscapes.

\textbf{BDA:} Outperforms BPSO and BGSA on 11 of 13 binary functions, with most differences strongly significant ($p < 0.002$).

\textbf{MODA:} Consistently better than NSGA-II on all five ZDT variants by a large margin; competitive with MOPSO, outperforming it on ZDT3 and the three-objective variant.
For the propeller design problem, MODA produced 61 well-distributed Pareto optimal solutions under strong constraints, demonstrating practical utility despite the absence of a known true Pareto front for verification.

Four behavioural diagnostics (search history, trajectory, average fitness, convergence curve) on selected functions confirm monotonic fitness improvement and eventual convergence, providing qualitative support for the algorithm's theoretical properties.

\bigskip

\section*{6.\ Assessment of conclusions}

The authors' primary claims --- that DA is a competitive single-objective optimiser and that BDA and MODA outperform their respective comparators --- are largely substantiated by the evidence.
The use of the Wilcoxon test and 30-run averaging follows sound methodology, and the breadth of benchmarks is commendable.

However, several caveats deserve attention.
First, the comparison set is narrow: PSO and GA are the only single-objective baselines.
Contemporary algorithms such as Differential Evolution or Grey Wolf Optimiser (also authored by Mirjalili) are absent, which limits the strength of the superiority claim.
Second, results on composite functions are mixed and the statistical advantage over PSO is modest, suggesting that DA's exploration/exploitation balance may not consistently generalise to the highly deceptive landscapes most characteristic of real-world problems.
Third, the propeller design comparison lacks an independently known true Pareto front, so MODA's performance cannot be benchmarked against an absolute standard.

For the purposes of this assignment, DA is an attractive candidate: its five operators and explicit food/enemy mechanism translate naturally to a continuous search space such as the trading-bot parameter optimisation in Part~2, and the algorithm is straightforward to implement from the pseudocode provided.
The publicly available MATLAB source code facilitates cross-verification.

\bibliographystyle{IEEEtran}
\begin{thebibliography}{99}

\bibitem{mirjalili2016dragonfly}
S.~Mirjalili, ``Dragonfly algorithm: a new meta-heuristic optimization technique for solving single-objective, discrete, and multi-objective problems,''
\textit{Neural Comput.\ Appl.}, vol.~27, no.~4, pp.~1053--1073, 2016.

\bibitem{kennedy1995pso}
J.~Kennedy and R.~Eberhart, ``Particle swarm optimization,'' in
\textit{Proc.\ IEEE Int.\ Conf.\ Neural Networks}, 1995, pp.~1942--1948.

\bibitem{wolpert1997nfl}
D.~H.~Wolpert and W.~G.~Macready, ``No free lunch theorems for optimization,''
\textit{IEEE Trans.\ Evol.\ Comput.}, vol.~1, no.~1, pp.~67--82, 1997.

\bibitem{reynolds1987flocks}
C.~W.~Reynolds, ``Flocks, herds and schools: a distributed behavioral model,''
\textit{ACM SIGGRAPH Comput.\ Graph.}, vol.~21, no.~4, pp.~25--34, 1987.

\bibitem{saremi2014transfer}
S.~Saremi, S.~Mirjalili, and A.~Lewis, ``How important is a transfer function in discrete heuristic algorithms,''
\textit{Neural Comput.\ Appl.}, pp.~1--16, 2014.

\bibitem{coello2004mopso}
C.~A.~C.~Coello, G.~T.~Pulido, and M.~S.~Lechuga, ``Handling multiple objectives with particle swarm optimization,''
\textit{IEEE Trans.\ Evol.\ Comput.}, vol.~8, no.~3, pp.~256--279, 2004.

\bibitem{cui2009boid}
Z.~Cui and Z.~Shi, ``Boid particle swarm optimisation,''
\textit{Int.\ J.\ Innov.\ Comput.\ Appl.}, vol.~2, pp.~77--85, 2009.

\bibitem{kadrovach2002particle}
B.~A.~Kadrovach and G.~B.~Lamont, ``A particle swarm model for swarm-based networked sensor systems,'' in
\textit{Proc.\ ACM Symp.\ Appl.\ Comput.}, 2002, pp.~918--924.

\bibitem{cui2009alignment}
Z.~Cui, ``Alignment particle swarm optimization,'' in
\textit{Proc.\ 8th IEEE Int.\ Conf.\ Cognitive Informatics}, 2009, pp.~497--501.

\end{thebibliography}

\end{document}
